% Generated draft: requires mathematical and source review.
\documentclass[11pt]{article}
\usepackage[margin=25mm]{geometry}
\usepackage{fontspec}
\setmainfont{Latin Modern Roman}
\usepackage{amsmath,amssymb,mathtools,unicode-math}
\setmathfont{Latin Modern Math}
\usepackage{enumitem,needspace,xurl,hyperref}
\setlength{\parindent}{0pt}
\setlength{\parskip}{5pt}
\setlength{\emergencystretch}{3em}
\newcommand{\N}{\mathbb{N}}
\newcommand{\Z}{\mathbb{Z}}
\newcommand{\Q}{\mathbb{Q}}
\newcommand{\R}{\mathbb{R}}
\newcommand{\C}{\mathbb{C}}
\newcommand{\F}{\mathbb{F}}
\newcommand{\A}{\mathbb{A}}
\newcommand{\PP}{\mathbb{P}}
\newcommand{\E}{\mathbb{E}}
\newcommand{\eps}{\varepsilon}
\newcommand{\dd}{\,\mathrm{d}}
\begin{document}
\section*{UnsolvedMath ID: 30000774}
\textbf{Source problem code:} catalogue entry
\subsection{Definitions and mathematical statement}
Let \(\Omega\subseteq\R^d\) be the spatial domain, \(u:\Omega\to\R^d\) the displacement, and \(z\) the scalar damage variable, with admissible spaces and constraints still to be specified from the report. Write \(e(u)=(\nabla u+\nabla u^{\mathsf T})/2\). The supplied complete-damage stored energy is
\[
 \widehat{\mathcal E}(u,z)=\int_\Omega
 \left[W(x,e(u),z)+\frac{\kappa}{r}|\nabla z|^r\right],\dd x,
\]
with the usual positive gradient coefficient \(\kappa\) and exponent \(r\geq1\). The source concerns an existence and convergence theorem for energetic solutions that is available for a quadratic dependence of \(W\) on strain and the supercritical range \(r>d\). It asks for two extensions: to suitable general nonquadratic strain-energy densities \(W\), and to the full lower range
\[
 1\leq r\leq d,
\]
including the critical exponent and the endpoint \(r=1\). The intended target includes the regime where both extensions are needed; they are not restrictions to one particular polynomial density or to \(r=d\) alone.

In an energetic formulation the stored energy is supplemented by external loading, admissibility conditions, and a dissipation distance \(\mathcal D\), possibly incorporating damage irreversibility. Schematically, for the full time-dependent energy \(\mathcal E\), global stability and energy balance take the form
\[
 \mathcal E(t,u(t),z(t))\leq
 \mathcal E(t,\widetilde u,\widetilde z)+\mathcal D(z(t),\widetilde z)
 \quad\text{for admissible competitors},
\]
\[
 \mathcal E(t,u(t),z(t))+\operatorname{Var}_{\mathcal D}(z;[0,t])
 =\mathcal E(0,u(0),z(0))+
   \int_0^t\partial_s\mathcal E(s,u(s),z(s))\,\dd s.
\]
Here \(\operatorname{Var}_{\mathcal D}\) denotes the supremum of sums of dissipation distances along partitions. These displays explain the standard energetic notion; the record does not specify the data needed to instantiate it or the approximation whose energetic limit is requested. Complete damage can remove elastic coercivity in a damaged region, so no uniformly positive elastic modulus has been inserted as a replacement assumption.

The requested extension is therefore conditional on recovering the original theorem's loading, dissipation, boundary, initial-data, admissibility, and approximation hypotheses and identifying appropriate nonquadratic assumptions. It is not an assertion for arbitrary functions \(W\) without growth, regularity, or lower-semicontinuity conditions.
\subsection{Short English statement}
The problem asks to extend an energetic existence and limit theorem for complete damage beyond quadratic elasticity and beyond the exponent condition \(r>d\), reaching every \(1\leq r\leq d\). The stored energy is given, but the remaining evolutionary model and approximation assumptions must be recovered before a precise theorem can be stated.
\subsection{Sources}
\begin{itemize}
\item Local UnsolvedMath registry, permanent ID \(30000774\), exact source code \texttt{catalogue entry}, proposed year 2007; no individual proposer is supplied.
\item \emph{Phase Transitions (2008)}, Oberwolfach Reports, DOI 10.4171/owr/2007/28, pp. 1575--1650. \url{https://doi.org/10.4171/owr/2007/28}.
\item The supplied \texttt{corrected\_verified} note reports direct checking of the complete-damage energy and both extensions of the quadratic, supercritical-exponent theorem. The open label and assessment checked 2026-08-21 state that the report leaves these extensions unresolved and that no general covering theorem was found in that triage. This is supplied dated assessment, not an independent search or status determination.
\end{itemize}
\end{document}
